% ---
% RESUMOS
% ---

% resumo em português
\setlength{\absparsep}{18pt} % ajusta o espaçamento dos parágrafos do resumo
\begin{resumo}
Este trabalho apresenta a implementação e validação do estimador de estados 
multissensorial MUSE (Multi-Sensor State Estimator) no robô quadrúpede Unitree 
Go2 (versão EDU). O objetivo central é a aplicação de uma estratégia de estimação
 robusta que combine dados proprioceptivos (IMU e encoders) e exteroceptivos 
 (LiDAR ou câmera) para garantir a estabilidade da locomoção em terrenos 
 complexos. Diferente das abordagens convencionais, o projeto utiliza o 
 middleware DLS2 (Dynamic Legged Systems) e foca na redução do drift de posição 
 por meio de módulos de detecção de deslizamento. A fundamentação teórica descreve 
 a cinemática do robô com 12 graus de liberdade, a arquitetura de observadores em 
 cascata (NLO e XKF) e a fusão sensorial. A validação inicial do estimador será 
 realizada em ambiente simulado no Gazebo, com análise de telemetria via 
 PlotJuggler e visualização de dados espaciais através da biblioteca PCL, 
 buscando preparar o sistema para aplicações em hardware real.

 \textbf{Palavras-chave}: Estimação de estado. Robô quadrúpede. Fusão de sensores. MUSE. Unitree Go2. Detecção de deslizamento.
\end{resumo}
